\subsubsection{Norma del conmutador de las proyecciones aritméticas}

Se obtiene el prefactor electrónico sin utilizar \(\alpha\), la torsión
o una unidad dimensional. El dominio es la tercera realización de APP,
\(\mathcal A_3=\{1,\ldots,9\}^3\), provista de medida uniforme.
El espacio de funciones es \(\ell^2(\mathcal A_3;\mathbb R)\), con producto
interno
\[
 \langle f,g\rangle=\frac1{729}\sum_{x\in\mathcal A_3}f(x)g(x).
\]
En este bloque, \([P,Q]=PQ-QP\) denota el conmutador aditivo de operadores.
Definimos
\[
 \Sigma(x)=\rho_9(x_1+x_2+x_3),\qquad
 \Pi(x)=\rho_9(x_1x_2x_3).
\]
Las esperanzas condicionales \(P_\Sigma\) y \(P_\Pi\) son las
proyecciones ortogonales sobre las funciones constantes en cada fibra
aditiva y multiplicativa, respectivamente. Así, para una función \(f\),
\[
 (P_\Sigma f)(x)
 =\frac{1}{|\Sigma^{-1}(\Sigma(x))|}
   \sum_{y:\,\Sigma(y)=\Sigma(x)}f(y),
\]
y análogamente para \(P_\Pi\). Las dos proyecciones actúan sobre el
mismo espacio; no representan dos soportes aritméticos independientes.

\begin{lemma}[Dos proyecciones y ángulos principales]
Si \(s\in[0,1]\) es un valor singular del emparejamiento entre bases
ortonormales de los rangos de dos proyecciones \(P,Q\), su bloque no
trivial contribuye a \(\|[P,Q]\|\) con
\(s\sqrt{1-s^2}\).
\end{lemma}
\begin{proof}
En el plano principal correspondiente se pueden escribir
\[
 P=\begin{pmatrix}1&0\\0&0\end{pmatrix},\qquad
 Q=\begin{pmatrix}
 s^2&s\sqrt{1-s^2}\\
 s\sqrt{1-s^2}&1-s^2
 \end{pmatrix}.
\]
Su conmutador es
\[
 [P,Q]=s\sqrt{1-s^2}
 \begin{pmatrix}0&1\\-1&0\end{pmatrix}.
\]
La matriz final es ortogonal, por lo que su norma es uno. Los bloques
comunes u ortogonales tienen conmutador nulo. La descomposición ortogonal
en planos principales da la afirmación y reduce la norma total al máximo
de estas contribuciones.
\end{proof}

\begin{proposition}[Norma exacta de incompatibilidad]
Sobre \(\ell^2(\mathcal A_3)\),
\begin{equation}
 \|[P_\Sigma,P_\Pi]\|=\frac{\sqrt3}{4}.
 \label{ivloc:norma}
\end{equation}
\end{proposition}
\begin{proof}
Sea \(N_{ab}=|\Sigma^{-1}(a)\cap\Pi^{-1}(b)|\). La enumeración entera
de los tres símbolos proporciona
\[
 N=9\begin{pmatrix}
 0&1&1&0&1&1&0&1&4\\
 1&0&1&1&0&1&1&0&4\\
 1&0&2&0&1&2&0&0&3\\
 0&1&1&0&1&1&0&1&4\\
 1&0&1&1&0&1&1&0&4\\
 0&0&2&1&0&2&0&1&3\\
 0&1&1&0&1&1&0&1&4\\
 1&0&1&1&0&1&1&0&4\\
 0&1&2&0&0&2&1&0&3
 \end{pmatrix}.
\]
Cada fila suma \(81\); las sumas de columna son
\[
 (m_1,\ldots,m_9)=(36,36,108,36,36,108,36,36,297).
\]
Las indicatrices normalizadas de las fibras tienen matriz de
emparejamiento \(C_{ab}=N_{ab}/\sqrt{81m_b}\). Por consiguiente,
\[
 M=CC^{\mathsf T},\qquad
 M_{ac}=\sum_{b=1}^9\frac{N_{ab}N_{cb}}{81m_b}
\]
es una matriz racional cuyo espectro consta de los cuadrados de los
valores singulares requeridos.

La multiplicación exacta muestra que \((1,\ldots,1)\),
\((1,-1,0)^3\) y \((1,1,-2)^3\) son autovectores de autovalores
\(1\), \(1/4\) y \(7/132\), respectivamente. El subespacio de
vectores soportados en \(\{3,6,9\}\) cuya suma es cero tiene dimensión
dos y autovalor \(1/18\). Los vectores de suma cero soportados en
\(\{1,4,7\}\) o en \(\{2,5,8\}\) forman un núcleo de dimensión
cuatro. Estos subespacios son independientes y sus dimensiones suman
nueve. Queda, pues, probado el espectro completo
\[
 \operatorname{spec}(M)=
 \left\{1,\frac14,\frac7{132},\frac1{18},\frac1{18},0,0,0,0\right\}.
\]
El máximo de \(\lambda(1-\lambda)\) en este conjunto es \(3/16\),
alcanzado en \(\lambda=1/4\). El lema anterior da
\(\sqrt{3/16}=\sqrt3/4\).
\end{proof}

El modo \((1,-1,0)^3\) es precisamente la distribución de valores del
selector trítico de fase en la coordenatización nonádica. La norma escalar retiene
una medida de incompatibilidad, pero no conserva por sí sola la
orientación de ese modo ni reconstruye los dos proyectores. Esa memoria
permanece en el estado de ruta que acompaña la lectura escalar.

\begin{proposition}[Álgebra del bloque electrónico local]
\label{ivloc:algebra}
En el plano principal que alcanza \eqref{ivloc:norma}, sean
\[
 P=\begin{pmatrix}1&0\\0&0\end{pmatrix},\qquad
 Q=\begin{pmatrix}1/4&\sqrt3/4\\\sqrt3/4&3/4\end{pmatrix},\qquad
 S=2P-I,\qquad J=\frac4{\sqrt3}[P,Q].
\]
Entonces \(S^2=I\), \(J^2=-I\), \(SJ=-JS\), y el álgebra real
generada por \(S,J\) es \(M_2(\mathbb R)\), realización de
\(\mathrm{Cl}_{1,1}\). Además,
\[
 n_\pm=\frac{S\pm J}{2},\qquad
 n_\pm^2=0,\qquad n_+n_-+n_-n_+=I.
\]
\end{proposition}
\begin{proof}
Se tiene \(S=\operatorname{diag}(1,-1)\) y
\(J=\bigl(\begin{smallmatrix}0&1\\-1&0\end{smallmatrix}\bigr)\).
La multiplicación verifica las tres relaciones. Las matrices
\(I,S,J,SJ\) son linealmente independientes y forman una base de
\(M_2(\mathbb R)\), lo que prueba la afirmación sobre el álgebra.
Finalmente, \((S\pm J)^2=S^2+J^2\pm(SJ+JS)=0\), y la suma de los
productos cruzados es \((S^2-J^2)/2=I\).
\end{proof}

En este bloque conviven una dirección elíptica, una hiperbólica y dos
direcciones nulas parabólicas. El resultado describe una representación
matricial local de la aritmética de las dos hojas. No identifica toda APP
con \(M_2(\mathbb R)\), ni con un cuadrado latino cuántico; tales
afirmaciones exigirían otros dominios y otras relaciones. Tampoco la
exponencial pertenece exclusivamente al régimen parabólico: puede
aplicarse a cualquiera de los generadores de este bloque.
